\chapter{Textures and Render Targets}
\label{ch:textures-rendertargets}

This chapter is the one in Part~\ref{part:graphics-machine} where the gap between ``the API exists'' and ``the API
does what it says on every renderer'' is widest. \cnaclass{Texture2D} is close to fully
faithful; \cnaclass{Texture3D} and \cnaclass{TextureCube} carried a real, load-bearing
structural deviation from FNA for a stretch of this project's history, closed since (see
\S\ref{sec:texture3d-inheritance}) but worth the full account of what changed and what did
not; and render targets --- once the sharpest per-renderer divergence in the whole Graphics
namespace, before a wave of later fixes closed most of it, covered in full below --- still
carry two genuinely open gaps of their own. All three are covered here together, with the
specific per-renderer facts intact rather than smoothed into ``mostly works.''

\section{Texture: a shared Color gate with one renderer-specific branch}

\cnaclass{Texture} (inheriting \cnaclass{GraphicsResource}) exposes \texttt{Format}
(\cnaclass{SurfaceFormat}) and \texttt{LevelCount}, plus a set of static helpers ported from
FNA's own internal size-calculation methods --- made \texttt{public static} in CNA rather
than FNA's \texttt{internal}, because three of the four real call sites in CNA are not even
\cnaclass{Texture} subclasses. The base \texttt{ValidateFormat(format)} method accepts only
\cnaclass{SurfaceFormat::Color}; every non-Skia build routes ordinary \cnaclass{Texture2D} and
\cnaclass{RenderTarget2D} construction through that gate, and \cnaclass{Texture3D} plus
\cnaclass{TextureCube} use it on every renderer.

The compile-time Skia product is the deliberate exception. Its \cnaclass{Texture2D} validator
admits 26 of the 27 enum values and retains each promoted native layout, including compressed
DXT/BC7 blocks; \texttt{Dxt5SrgbEXT} is the one refusal. Its render-target validator admits the
14 formats that FNA reports renderable. The
exception is therefore specific to two 2D resource classes, not permission to infer broad format
support from the enum alone. Separately, the portable \texttt{Texture2D::FromStream} DDS route
recognizes DXT1/3/5 and CPU-decompresses them to RGBA8 before ordinary upload.

\section{A structural deviation that closed since this chapter's own first pass}
\label{sec:texture3d-inheritance}

Real XNA's \texttt{Texture3D} and \texttt{TextureCube} both inherit \texttt{Texture}, which
is what lets any of them sit in \texttt{GraphicsDevice.Textures[slot]} and be bound to a
custom shader like an ordinary 2D texture. An earlier pass of this chapter reported that
CNA's \cnaclass{Texture3D} and \cnaclass{TextureCube} did \textbf{not} inherit
\cnaclass{Texture} at all, tracked as Task~863 --- reading the real headers directly today
shows that claim is stale: \texttt{Texture3D.hpp} and \texttt{TextureCube.hpp} both now
declare \texttt{class Texture3D : public Texture} / \texttt{class TextureCube : public
Texture}, matching real XNA exactly. \texttt{Texture3DTests.cpp}'s own regression coverage
names precisely what this fix restores --- assignment into a \cnaclass{TextureCollection}
(\texttt{GraphicsDevice.Textures[slot] = texture3D}), a real compile-time impossibility
before Task~863 closed, plus a correctness follow-on: \texttt{Texture3D::Dispose(bool)} now
calls through to \texttt{Texture::Dispose(bool)} after releasing its own renderer handle, so
disposing a bound \cnaclass{Texture3D} correctly unbinds it from
\texttt{GraphicsDevice.Textures} / \texttt{VertexTextures}, the same unbind-on-dispose contract
\cnaclass{Texture2D} already honored and \cnaclass{Texture3D} previously skipped entirely.

The practical consequence this earlier pass described --- ``a custom
\cnaclass{ShaderEffect} has no way to bind a \cnaclass{Texture3D} or \cnaclass{TextureCube}''
--- is also closed, and by a more direct route than the generic texture-slot mechanism the
class-hierarchy fix alone would suggest: \cnaclass{ShaderEffect} gained dedicated
\texttt{SetTexture(int, TextureCube\&)} (Task~1081) and \texttt{SetTexture(int, Texture3D\&)}
(Task~863) overloads, each binding straight to the matching GL texture target
(\texttt{glBindTexture(GL\_TEXTURE\_CUBE\_MAP, ...)} / \texttt{GL\_TEXTURE\_3D}) rather than
routing through the ordinary 2D texture-slot path. This is EasyGL-only as of this writing ---
the other nine renderers were explicitly deferred, per the same task's own scope note --- so
``a custom shader can sample a volume or cube texture'' is currently a real, tested EasyGL
capability, not yet a cross-renderer one. \cnaclass{EnvironmentMapEffect}'s own
\texttt{EnvironmentMap} field, mentioned in an earlier pass as the \emph{only} way any effect
could reach a cube texture, is now one of two paths rather than the sole exception.

\section{Texture2D}

\cnaclass{Texture2D} is copyable and movable (its renderer is a \texttt{shared\_ptr}-held
\texttt{ITextureRenderer}), constructible from raw dimensions, from a file path (two
\texttt{CNAEXT} overloads --- real XNA has no such constructor, since real XNA loads textures
exclusively through the compiled content pipeline), or via \texttt{FromStream} (device plus
stream, or a five-argument form that resizes/crops to an explicit width, height, and
\texttt{zoom} fit-vs-cover flag). \texttt{SetData} / \texttt{GetData} cover full-array and
partial (level, rectangle, start index, element count) forms, all \cnaclass{Color}-only.
\texttt{SaveAsPng} / \texttt{SaveAsJpeg} round out the surface, alongside several
\texttt{CNAEXT} escape hatches (\texttt{GetRenderer()}, \texttt{GetCpuPixelsWeak()},
\texttt{CreateFromPixels}) that exist for interop with CNA's own content-pipeline and testing
code rather than for game code.

\subsection{A worked example: building a texture from raw pixel data}

Every one of the API traits described above is visible in one small, real pattern: building a
procedural texture entirely from \cnaclass{Color} data, with no file on disk at all --- a
checkerboard, grounded directly in the real constructor and \texttt{SetData} / \texttt{GetData}
signatures in \texttt{Texture2D.hpp}:

\begin{lstlisting}
const int size = 64;
Texture2D checker(getGraphicsDeviceProperty(), size, size); // defaults to SurfaceFormat::Color

std::vector<Color> pixels(size * size);
for (int y = 0; y < size; ++y) {
    for (int x = 0; x < size; ++x) {
        bool light = ((x / 8) + (y / 8)) % 2 == 0;
        pixels[y * size + x] = light ? Color::White : Color::Black;
    }
}
checker.SetData(pixels.data(), static_cast<int>(pixels.size()));

// Read the CPU shadow back (this does not prove what the GPU samples):
Color topLeft;
checker.GetData(&topLeft, 0, 1);
\end{lstlisting}

Two constraints from this chapter's opening section apply directly to this example. The
two-argument \texttt{Texture2D(device, width, height)} constructor defaults to
\cnaclass{SurfaceFormat::Color}, so \cnaclass{Color} is the portable path across renderer
families. On a non-Skia build, passing any of the other 26 \cnaclass{SurfaceFormat} values to
the explicit-format constructor throws before upload. The public class does now provide typed
\texttt{SetData}/\texttt{GetData} overloads for packed vectors, floats, vectors, half formats,
bytes, and unsigned shorts; those routes become materially useful with Skia's promoted native
layouts. They are concrete overloads rather than XNA's arbitrary \texttt{T[]} template.

\subsection{Two SetData paths, and why GetData can hide a failed upload}
\label{sec:texture2d-update-contract}

The simple whole-texture call in the example and the detailed level/rectangle overload do
not update the resource the same way. \texttt{SetData(Color*, count)} converts the complete
level~0 image, replaces this object's shared renderer reference, and calls
\texttt{CreateTexture()} for a replacement resource. The detailed overload instead patches
CNA's CPU shadow first. At level~0 it sends the \emph{whole reconstructed level} through
\texttt{ITextureRenderer::UpdatePixels()}, even when the caller supplied a small rectangle;
above level~0 it sends that level's complete shadow through
\texttt{UpdatePixelsLevel()}. This differs from FNA, whose detailed overload forwards the
requested rectangle directly to \texttt{FNA3D\_SetTextureData2D()}.

That CPU shadow changes what a test proves. An ordinary \cnaclass{Texture2D::GetData()} reads
\texttt{cpuPixels\_} or \texttt{extraMipLevels\_}; it does not read the sampled GPU resource.
The source file for \path{easygl_texture2d_mip_test.cpp} says this explicitly: its three-level
``round trip'' is pure CPU-shadow verification with no framebuffer readback. It would still
return every color supplied by \texttt{SetData()} if the corresponding renderer update were
an empty body. GPU sampling, a native renderer readback, or a rendered pixel is required to
establish the other half.

That distinction exposes a live split. EasyGL, Vulkan, BGFX, WebGPU, and D3D9/11/12 allocate
and update real \cnaclass{Texture2D} mip subresources. SDL\_RENDERER, ASCII, Canvas, and FREEDIRECT
reject a level-above-zero upload with \texttt{std::runtime\_error}. Software explicitly
discards it, Headless records the call without storing a renderer mip, and SDL GPU inherits
the interface's empty \texttt{UpdatePixelsLevel()} body while allocating only one physical
level even when public \texttt{LevelCount} reports a full chain. In all three discard cases,
the public CPU shadow still remembers the supplied level, so a following
\texttt{GetData(level,...)} can look successful while rendering can never consume it.

WebGPU has the opposite extra behavior: writing level~0 automatically regenerates every
higher level with a real linear-filtered render-pass cascade. That prevents undefined mip
content, but differs from FNA and means a later level~0 update overwrites any authored
higher levels written earlier. Chapter~\ref{ch:backend-architecture},
\S\ref{sec:texture2d-update-matrix}, gives the complete renderer and evidence matrix.

\subsection{FromStream format support}
\label{sec:texture-streams}

Format detection in \texttt{FromStream} starts with a minimal internal DDS-header sniffer
(recognizing only the DXT1/3/5 FourCC codes, decoding via an internal \texttt{DxtUtil}), and
falls through to SDL3\_image's \texttt{IMG\_Load\_IO} --- which auto-detects the container
format from the byte stream itself --- for everything else. Verified round-trip formats:
PNG, JPEG (lossy, within tolerance), 24-bit uncompressed BMP, and DDS (DXT1/3/5). AVIF, TIFF,
and WebP are present in the linked SDL3\_image build's own dependencies but are not
independently tested by CNA's own suite. The \texttt{zoom} parameter on the five-argument
overload matters more than its name suggests: \texttt{zoom=false} uses a simplified
largest-dimension-fit heuristic (not a general \texttt{min(w/w0, h/h0)} bounding-box fit --- it
assumes a roughly square target box), while \texttt{zoom=true} scales up and center-crops to
exactly fill the requested box.

\section{Texture3D and TextureCube}

Beyond the structural deviation described above, \cnaclass{Texture3D} and
\cnaclass{TextureCube} are both non-copyable (each owns a unique GPU handle) but movable ---
\cnaclass{Texture3D} specifically gained move construction/assignment to satisfy a glTF
content-pipeline reader's needs. Their requested \cnaclass{SurfaceFormat} is not silently
ignored: the shared \texttt{Texture::ValidateFormat()} gate rejects every value except
\cnaclass{SurfaceFormat::Color} before a renderer is created, so every surviving resource is
RGBA8 by construction.

Neither type keeps a \cnaclass{Texture2D}-style CPU pixel shadow. Every valid
\texttt{GetData()} on a constructed renderer therefore reaches that renderer's own readback
path. The old support document's claim that only EasyGL really implements this is now stale:
EasyGL, Vulkan, BGFX, WebGPU, SDL GPU, and all three Direct3D renderers have real GPU readback
for both resource types. Their mechanisms differ --- framebuffer reads, staging copies,
native lock/map calls, or a temporary BGFX transfer texture --- but all eight return actual
sub-region bytes. Software adds a real level-0 CPU implementation for
\cnaclass{TextureCube}, while its \cnaclass{Texture3D} factory is absent. Headless deliberately
returns transparent-black placeholders. Chapter~\ref{ch:backend-architecture}'s
\S\ref{sec:texture-volume-readback-matrix} gives the full factory/readback matrix, including
the separate and much narrower \cnaclass{RenderTargetCube} result.

Authored mip-level upload and readback are likewise real on both ordinary resource types
across those eight GPU renderers. Exact public mip round trips are registered for EasyGL,
Vulkan, BGFX, WebGPU, and SDL GPU; the Direct3D implementations allocate and address mip
subresources too, although their cited smoke-test readbacks concentrate on discriminating
level-0 subregions. This statement is about explicitly supplied mip data, not a promise that
every renderer automatically generates a mip chain after a level-0 upload.
\texttt{CubeMapFace} validation, by contrast, is a place CNA is \emph{stricter} than FNA: an
out-of-range face throws \texttt{std::out\_of\_range} at the public API layer, a safety check
real XNA never performs at all.

\subsection{A worked example: six faces, six SetData calls}

\cnaclass{TextureCube}'s constructor takes a single face size (every face is square and
identical in size, per the real XNA contract) plus \texttt{mipMap} and \texttt{format} --- and,
per this chapter's own opening finding, \texttt{format} must be
\cnaclass{SurfaceFormat::Color} or construction throws immediately. Populating a solid-color
test cube (a stand-in for a real skybox's six photographic faces) is a direct, mechanical loop
over \cnaclass{CubeMapFace}'s six enumerators, since \texttt{SetData} takes the face as its
first argument:

\begin{lstlisting}
const int faceSize = 128;
TextureCube skybox(getGraphicsDeviceProperty(), faceSize, /*mipMap=*/false,
                    SurfaceFormat::Color);

const CubeMapFace faces[6] = {
    CubeMapFace::PositiveX, CubeMapFace::NegativeX,
    CubeMapFace::PositiveY, CubeMapFace::NegativeY,
    CubeMapFace::PositiveZ, CubeMapFace::NegativeZ,
};
const Color faceColors[6] = {
    Color::Red,    Color(128, 0, 0, 255),   // +X, -X
    Color::Green,  Color(0, 128, 0, 255),   // +Y, -Y
    Color::Blue,   Color(0, 0, 128, 255),   // +Z, -Z
};

std::vector<Color> pixels(faceSize * faceSize);
for (int i = 0; i < 6; ++i) {
    std::fill(pixels.begin(), pixels.end(), faceColors[i]);
    skybox.SetData(faces[i], pixels.data(), static_cast<int>(pixels.size()));
}
\end{lstlisting}

This loop can now be verified with a real \texttt{GetData()} round trip on any of the eight
GPU renderers above, rather than only on EasyGL as the stale support document still says.
The strongest reusable checks deliberately distinguish slices, faces, off-origin regions,
and mip levels; they do not merely assert that the call returned. Software can make the same
check at level~0 for a cube. \texttt{skybox} \emph{can} also be bound to a custom
\cnaclass{ShaderEffect} on EasyGL, per
\S\ref{sec:texture3d-inheritance} above --- either through \texttt{SetTexture(int,
TextureCube\&)} directly, or, since \cnaclass{TextureCube} now inherits \cnaclass{Texture},
through the ordinary \texttt{GraphicsDevice.Textures[slot]} assignment every 2D texture in
this chapter already uses.

\subsection{A worked example: sampling a volume texture from a custom shader}
\label{sec:texture3d-shader-example}

\texttt{modules/renderers/easygl/examples/\allowbreak{}easygl\_shadereffect\_texture3d\_test.cpp} is the real, EasyGL-only test
that proves \texttt{ShaderEffect::SetTexture(int, Texture3D\&)} genuinely reaches a
\texttt{sampler3D} uniform, not just that the call compiles. The volume is deliberately
tiny --- a $1\times1\times2$ texture, one texel per Z-slice --- so the whole capability proof
reduces to two texel reads instead of a real ray-marching setup:

\begin{lstlisting}
Texture3D volumeTex(device, 1, 1, 2, /*mipMap=*/false, SurfaceFormat::Color);
const Color slices[2] = {
    Color(255, 0, 0, 255), // Z=0
    Color(0, 0, 255, 255), // Z=1
};
volumeTex.SetData(slices, 2);

auto* fx = dynamic_cast<ShaderEffect*>(fxBase.get());
fx->Apply();
fx->SetTexture(0, volumeTex);
fx->SetUniformInt("VolumeSampler", 0);
fx->SetUniformVec3("coord", 0.5f, 0.5f, 0.25f); // samples Z-slice 0's texel centre
\end{lstlisting}

Sampling at $Z=0.25$ and $Z=0.75$ --- each slice's own texel centre, per
$(\text{sliceIndex} + 0.5) / \text{depth}$ --- reads back \texttt{(255,0,0,255)} and
\texttt{(0,0,255,255)} respectively, even with GL\_LINEAR filtering on both axes: linear
interpolation sampled exactly at a texel centre puts its full weight on that one texel, so
there is no blending artifact to account for. The real test's own false-positive guard is
worth naming as a technique in its own right, since it recurs throughout this book's worked
examples: a second, all-black ``decoy'' \cnaclass{Texture3D} is created and uploaded
\emph{after} the real one, so an unmutated \texttt{SetTexture()} call would leave GL's
texture-unit-0/\texttt{GL\_TEXTURE\_3D} binding pointing at the decoy rather than the
intended volume --- the two-slice colour check only passes for the right reason (a genuine,
per-call rebind) if disabling the \texttt{SetTexture()} call would make it fail.

\subsection{A worked example: rendering into one RenderTargetCube face}

\cnaclass{GraphicsDevice} has a dedicated overload for exactly this purpose ---
\texttt{SetRenderTarget(RenderTargetCube*, CubeMapFace)} --- distinct from the plain
\texttt{SetRenderTarget(RenderTarget2D*)} form used for ordinary 2D targets, and from the
\texttt{SetRenderTargets(const std::vector<RenderTargetBinding>\&)} form used for
simultaneous multi-target binding (\cnaclass{RenderTargetBinding}'s own
\texttt{(Texture*, CubeMapFace)} constructor exists specifically to build entries for that
vector form, not for single-target binding). Rendering a small scene into a single cube face
and then sampling that cube in a later draw is the pattern a real-time reflection probe or a
baked-environment-map tool would use:

\begin{lstlisting}
RenderTargetCube probe(getGraphicsDeviceProperty(), 256, /*mipMap=*/false,
                        SurfaceFormat::Color, DepthFormat::Depth24Stencil8);

device.SetRenderTarget(&probe, CubeMapFace::PositiveZ);
device.Clear(Color::CornflowerBlue);
// ... draw the scene as seen looking down +Z from the probe's location ...
device.SetRenderTarget(nullptr); // back to the backbuffer
\end{lstlisting}

Sampling and CPU readback are separate capabilities here. The code above can feed the cube to
an effect without ever calling \texttt{probe.GetData()}; nevertheless public rendered-face
readback is now exact on EasyGL, Vulkan, BGFX, WebGPU, SDL GPU, and D3D9/11/12. Skia provides a
separate exact six-surface CPU emulation. Unsupported identities now throw instead of converting
a renderer no-op into fabricated transparent black. Mip/MSAA scope still differs --- notably
D3D9 has only level~0 and WebGPU refuses a mipmapped cube target --- so the central matrix in
\S\ref{sec:texture-volume-readback-matrix} remains the authoritative per-level comparison.

The inherited \texttt{SetData} family has narrower renderer support still. It is legal at
compile time because \cnaclass{RenderTargetCube} inherits the whole
\cnaclass{TextureCube} upload family, exactly as it does in FNA. EasyGL implements the GPU upload
with its documented row convention, and Skia implements exact CPU-surface upload. Every other
constructed target returns failure from the renderer contract; the public call validates first
and then throws \cnaclass{NotSupportedException} instead of returning normally after changing
nothing. Thus the feature remains narrow, but the former silent compatibility gap is closed.

The unbind line now routes correctly on Vulkan, WebGPU, SDL GPU, BGFX, D3D9, D3D11 and D3D12.
The Direct3D~11/12 renderers explicitly track the active cube and finalize it before every
cube/2D/MRT/backbuffer transition; shared public readback, usage and per-face-MSAA fixtures make
that edge observable. EasyGL retains a narrower exception: switching directly between faces of
the same cube skips outgoing-face finalization. Per-face MSAA storage now preserves each face's
samples, but only a finalized face is resolved/generated for later sampling; insert a
backbuffer unbind between freshly rendered faces when all must immediately be sampled. The full
transition matrix and BGFX's renderer-local (not normalized public) cube-null no-op are recorded
in \S\ref{sec:rendertargetcube-default-contracts}.

Whether this actually produces a correct image is, per \S\ref{sec:rendertargets} below,
renderer-dependent in ways the call site cannot infer. EasyGL renders it correctly. Vulkan's
two formerly-tracked black-render defects (Tasks~875/876) are now both closed: clear-only
targets enter the used-target list, and the original cube-sampling regression test passes
repeatedly end to end, although the exact incidental change that retired the second defect was
not bisected. BGFX's former unconditional \texttt{static\_cast} type-confusion bug is likewise
fixed by shared cube-sampling interfaces and virtual handle dispatch. The remaining transition,
same-frame, and upload qualifications above still apply; ``Vulkan cube sampling is broken'' no
longer does.

\section[Render-target divergences]{Render targets: a once-sharp per-renderer divergence,
mostly closed since}
\label{sec:rendertargets}

\cnaclass{RenderTarget2D} inherits both \cnaclass{Texture2D} and \cnaclass{IRenderTarget},
which is what lets it be both drawn into and sampled as an ordinary texture afterward; its
\texttt{Dispose(bool)} correctly throws \texttt{InvalidOperationException} if it is still
bound to the device, matching FNA. An earlier explanation for
\cnaclass{RenderTargetCube}'s missing equivalent guard is now stale:
\cnaclass{RenderTargetCube} does inherit \cnaclass{TextureCube} and
\cnaclass{IRenderTarget}, while \cnaclass{TextureCube} itself inherits \cnaclass{Texture}, so
\cnaclass{RenderTargetBinding} can store it through a \texttt{Texture*}. The guard is still
absent for a different reason. The singular cube-target setter clears
\texttt{currentRenderTargets\_} and retains only a boolean ``some target is bound,'' not the
cube pointer; \cnaclass{RenderTargetCube} supplies no \texttt{Dispose(bool)} override of its
own to compensate. The compatibility gap remains, but the former class-hierarchy explanation
does not.

\subsection{\cnaclass{RenderTargetUsage}: three values, more than three behaviors}
\label{sec:rendertarget-usage-contract}

\cnaclass{RenderTargetUsage} has the expected three ordinals:
\texttt{DiscardContents}, \texttt{PreserveContents}, and \texttt{PlatformContents}.
\texttt{DiscardContents} says old attachment data need not survive a bind; it does
\emph{not} promise a readable replacement color. \texttt{PreserveContents} requests survival
even at a performance or memory cost. \texttt{PlatformContents} permits survival only when
cheap. In current FNA, both non-discard values are passed to the native layer as a preserve
request, while the exact enum still controls whether the managed layer issues its discard
clear.

CNA now centralizes those decisions in
\texttt{RenderTargetUsagePreservesContentsEXT()}: only Discard maps to false, while Preserve
and Platform both map to true, matching FNA. \cnaclass{RenderTarget2D} and
\cnaclass{RenderTargetCube} pass that Boolean through their renderer factories, and the shared
normalized bind route emits the deterministic black/max-depth/zero-stencil clear only for
Discard. An explicit \texttt{Clear()} remains a distinct ordered operation and wins over the
usage policy aspect by aspect.

Cube targets now use the same binding path rather than storing only an enum beside an
untracked face. A \texttt{RenderTargetBindingDescriptor} preserves the selected face,
dimensions, applied sample count, and renderer pointer; usage reaches the cube factory, and
the shared discard clear can inspect the cube's real depth/stencil attachments. The remaining
differences are native mechanisms and renderer capability --- for example WebGPU rejects MRT
--- not a missing cube usage parameter.

The backbuffer's corresponding
\texttt{PresentationParameters.RenderTargetUsage} is currently stored, cloned, and exposed
only. Returning to the backbuffer never consults it. FNA, by contrast, supplies it to native
target binding and uses it to choose the unbind-to-backbuffer discard clear. The first 2D
target's usage controls CNA's MRT implicit clear, which matches FNA's first-target policy in
shape, but deferred Vulkan MRT passes discard regardless and the public plural cube form
cannot reach a cube renderer at all.

Current evidence is substantially broader than the early non-MSAA 2D tests. Shared cube-usage,
per-face MSAA preservation, depth/stencil usage, pass-boundary, and plural-cube fixtures are
registered across capable renderers; Platform is asserted as preservation rather than left
implicit. Backbuffer usage and redundant-binding elision remain separate gaps: returning to
the backbuffer still ignores its presentation usage, and CNA still performs native/state work
for an identical repeated binding instead of FNA's early return. The complete source/evidence matrix is
\S\ref{sec:backend-rendertarget-usage-matrix}; the binding-identity consequences are
\S\ref{sec:graphicsdevice-target-binding-contract}.

The render-target support ledger once listed five further per-renderer divergences here
(Tasks~877--881) as open gaps, and an earlier pass of this chapter reproduced that list
faithfully. Reading every renderer's own current \texttt{.cpp} directly instead of trusting that
document's status column finds four of the five already closed --- each superseded by a later
task (903, 907, or 911) whose own fix comments the tracking document
was never updated to reflect, the identical ``a status that looked settled turned out to be
stale'' pattern Chapters~\ref{ch:easygl-backend} and~\ref{ch:vulkan-backend} both ran into
independently while auditing this exact same document for their own renderers. What is actually
still open, and what is now fixed:

\begin{description}
  \item[Cube sampling.] Sampling a \cnaclass{RenderTargetCube} face after rendering into it,
        then unbinding, works correctly on \textbf{EasyGL}. \textbf{Vulkan} now works in both
        discriminating
        regressions --- Task~875's clear-only case explicitly records the target, while
        Task~876's original render-all-faces/then-sample test no longer reproduces the black
        result (the exact fixing commit was not retroactively isolated).
        \textbf{BGFX}: fixed --- the unconditional \texttt{static\_cast} that once read a
        framebuffer-pool handle where a texture-pool handle was expected
        (Chapter~\ref{ch:bgfx-backend} names the exact mechanism) is gone, replaced by a real
        \texttt{dynamic\_cast} to a shared cube-samplable interface for the
        \cnaclass{EnvironmentMapEffect} path and virtual width/height dispatch for the ordinary
        \cnaclass{SpriteBatch} path.
  \item[Depth/stencil format fidelity.] Now real on all three renderers: each maps the actual
        requested \cnaclass{DepthFormat} ordinal to the correct native format (GL internal
        format plus attachment point on EasyGL; a per-instance Vulkan format; bgfx's own
        format enum), rather than the single, always-the-same format every renderer once
        allocated regardless of what was asked for.
  \item[Mipmap generation.] \textbf{EasyGL} already worked --- real per-level GPU storage,
        auto-regenerated via \texttt{glGenerateMipmap} on unbind, matching FNA3D's own real
        mechanism. \textbf{Vulkan} and \textbf{BGFX} now do too: a real, per-level mip-chain
        regeneration cascade (\texttt{vkCmdBlitImage} on Vulkan; bgfx's own equivalent) runs
        against a render target's just-rendered content once its render pass ends, rather than
        silently accepting and discarding the request the way both once did.
  \item[Multisample antialiasing.] \textbf{EasyGL} already worked: a real multisampled
        renderbuffer with a \texttt{glBlitFramebuffer} resolve on unbind. \textbf{Vulkan} and
        \textbf{BGFX} now genuinely implement it too, rather than honestly reporting
        \texttt{MultiSampleCount = 0} the way both once correctly did while the feature was
        still unimplemented.
\end{description}

Two further, project-wide findings from the same original list are also now closed, not open.
Switching the active render target correctly resets the shared public \texttt{Viewport} and
\texttt{Scissor\allowbreak{}Rectangle} values to the new target's (or the backbuffer's) own
dimensions on every renderer; whether those setters then change native rasterization is the separate
current renderer contract in \S\ref{sec:backend-viewport-scissor-matrix}.
\texttt{GraphicsDevice.Viewport} itself now has real GPU wiring on all three
hardware renderers investigated --- EasyGL applies a custom sub-region \cnaclass{Viewport}
on its ordinary 3D paths whether or not a render target is currently bound (SpriteBatch is the
documented exception), while Vulkan's own version
of the identical fix is honestly narrower, backbuffer-pass-only, for an architectural reason
Chapters~\ref{ch:vulkan-backend} and~\ref{ch:easygl-backend} both cover in full. And
\texttt{SetRenderTargets()} now genuinely enforces FNA's real four-target cap at the shared,
cross-renderer level --- throwing if a caller passes more than four bindings, regardless of
what any individual renderer's own internal buffer happens to be sized for --- closing what
this section once described as no renderer matching that real XNA limit at all.

\section{SurfaceFormat, without flattening the Skia exception}

The 27-value \cnaclass{SurfaceFormat} enum describes a wider contract than most configured
products implement. For 41 renderer families, ordinary texture construction still accepts only
\cnaclass{SurfaceFormat::Color}; compressed content loading commonly becomes RGBA8 before the GPU.
Skia is the bounded exception: 26 formats have native \cnaclass{Texture2D} storage and typed
transfer routes (all except \texttt{Dxt5SrgbEXT}), while only 14 are admitted for
\cnaclass{RenderTarget2D}. Volume and cube
textures remain Color-only even there. A format enumerator, a transferable CPU representation,
and render-target capability are therefore three separate claims.

One real, historical
bug is worth naming because of what it reveals about how subtle a graphics-correctness bug
can be: Vulkan's texture path once created images as \texttt{VK\_FORMAT\_R8G8B8A8\_SRGB}
(wrong --- \cnaclass{SurfaceFormat::Color} is linear, not sRGB) while the swapchain
\emph{separately} preferred a matching sRGB format. The two wrong encode/decode steps
approximately canceled out for ordinary textured content, which is exactly why the bug went
unnoticed for a long time --- it only became visible on non-textured content (flat vertex
colors, raw lighting output), where a mid-gray value of 128 was read back as 188. Both were
corrected to their \texttt{\_UNORM} equivalents.

\section{Buffers: VertexBuffer and IndexBuffer}

\cnaclass{VertexBuffer} and \cnaclass{IndexBuffer} both inherit \cnaclass{GraphicsResource}
and both maintain a private \texttt{cpuShadow\_} byte buffer purely to support
\texttt{GetData()} without requiring a real per-renderer GPU readback path --- since nothing
in CNA's own pipeline ever writes GPU-side data back into one of these buffers, a CPU-side
shadow copy is a faithful, sufficient implementation of read-your-own-writes semantics.
\cnaclass{VertexBuffer::SetData} / \texttt{GetData} are typed per concrete vertex struct ---
\cnaclass{VertexPositionColor},
\cnaclass{VertexPositionColorTexture},
\cnaclass{VertexPositionNormalTexture}, and
\cnaclass{VertexPositionTexture} --- plus a
\texttt{CNAEXT} skinned variant and a raw \texttt{SetDataRaw(data, count, stride)} escape hatch
for layouts with no dedicated XNA struct. \cnaclass{DynamicVertexBuffer} and
\cnaclass{DynamicIndexBuffer} add \texttt{IsContentLost} (always \texttt{false}) and
\texttt{SetData(..., SetDataOptions)} overloads. One half of the chapter's old simplification
is still true: the protected \cnaclass{VertexBuffer} and \cnaclass{IndexBuffer} constructors
receive a \texttt{dynamic} boolean from the two derived classes but deliberately leave the
parameter unnamed and never pass it to a renderer factory. Choosing \texttt{Dynamic*} therefore
does not itself select different storage. D3D9 and D3D11 allocate their ordinary and dynamic
buffers through the same always-dynamic native path; BGFX likewise uses dynamic handles for
both. The \emph{per-upload} hint is no longer universally ignored, however:

\begin{itemize}
  \item EasyGL maps \texttt{Discard} to buffer orphaning plus a new sub-data upload and
    \texttt{NoOverwrite} to an in-place sub-data update. Its vertex path makes this distinction
    whenever storage already exists. Its index path has one surprising coupling: when context
    recovery is disabled, no renderer CPU shadow is retained, so its
    \texttt{NoOverwrite}-eligibility check fails and falls back to a fresh data allocation.
  \item SDL GPU's native \texttt{cycle} flag is true for \texttt{None} and
    \texttt{Discard}, but false for \texttt{NoOverwrite}. D3D9 uses the corresponding
    \texttt{DISCARD}/%
    \texttt{NO}\allowbreak{}\texttt{OVERWRITE} lock flags once storage exists (a fresh
    allocation forces discard). D3D11 uses the corresponding
    \texttt{WRITE\_}\allowbreak{}\texttt{DISCARD}/%
    \texttt{WRITE\_}\allowbreak{}\texttt{NO\_OVERWRITE} map modes.
  \item Vulkan and BGFX inherit the interface default, which discards the option and invokes
    plain \texttt{SetData}. WebGPU and D3D12 explicitly receive but ignore it before their
    queue-write or synchronous staging-copy path. Software and Headless likewise replace their
    CPU/trace data without interpreting the hint. The four 2D-only renderers reject buffer
    construction before any of these methods can be reached.
\end{itemize}

The real \texttt{easygl\_dynamic\_buffer\_stress\_test.cpp} exercises the distinction through
the public classes rather than calling a renderer directly. Across twelve frames it cycles
\texttt{None}, \texttt{Discard}, and \texttt{NoOverwrite}, uploads a differently colored
six-vertex quad and six indices, draws, then reads the center pixel and checks the two public
capacities. That is a genuine pixel proof that all three routes preserve current data through
repeated use. It is not, by itself, proof of which native allocation primitive ran; the
orphan/sub-data distinction is established by the implementation read above.

\subsection{A shared readback hole in the options-taking path}

That audit exposes a separate bug above every renderer. The ordinary typed
\path{VertexBuffer::SetData} and \path{IndexBuffer::SetData} overloads update the shared
\texttt{cpuShadow\_} after uploading. None of the protected
\path{SetDataWithOptions} overloads does. A \cnaclass{DynamicVertexBuffer} or
\cnaclass{DynamicIndexBuffer} created with \texttt{BufferUsage::None} can therefore inherit
\texttt{GetData}, but after its normal options-taking \texttt{SetData} call the shared shadow
is still empty (or stale after an earlier ordinary upload): a fresh buffer throws
\path{ArgumentOutOfRangeException} instead of returning the bytes just written.
\path{BufferUsage::WriteOnly} still throws \path{NotSupportedException} by design; that
documented restriction does not excuse the broken non-write-only case. This is shared public
code, so changing graphics renderer cannot avoid it.

\section{Vertex layout: fixed by stride, except for EasyGL raw declarations}
\label{sec:vertex-layout}

The ordinary typed paths still choose fixed GPU layouts from the uploaded byte stride. The
current project-wide family has grown from five to eight shapes: 16, 20, 24, and 32 bytes for
the four stock XNA vertex types; 48 for CNA's PBR tangent layout; 52 for GPU skinning; 56 for
the same skinned data plus vertex color; and 68 for tangent-space PBR plus skinning. Exact
shader/effect combinations remain renderer-specific, so recognizing a stride is not a promise
that every effect can consume it.

There is now one deliberate exception to the stride-only rule.
\path{VertexBuffer::SetDataRaw()} pushes the buffer's own
\cnaclass{VertexDeclaration} element list through
\path{IVertexBufferRenderer::SetVertexDeclaration()} immediately before upload. That
interface method defaults to a no-op, and EasyGL is currently its only override. EasyGL maps
all twelve XNA element-format values to the correct GL attribute shape and uses the
declaration's element order as GLSL locations \(0,1,\ldots\); the semantic
\cnaclass{VertexElementUsage} does not choose the location.

The pixel proof is stronger today than its original ``non-standard stride'' description.
\path{easygl_shadereffect_custom_vertex_layout_test.cpp} supplies five elements in a
48-byte record: three \texttt{Vector3} fields, one \texttt{Vector2}, and a trailing normalized
\texttt{Color}. Forty-eight bytes is now also the size of CNA's fixed PBR layout, but that
layout uses a four-float tangent and has no trailing color. The test encodes normal, tangent,
UV, and color inputs into the output pixel; its two distinct, tolerance-bounded readbacks
therefore prove that EasyGL followed the declaration's offsets rather than merely selecting
its different fixed 48-byte case.

The practical rule is consequently renderer-specific. An arbitrary
\cnaclass{VertexDeclaration} is a real EasyGL+\cnaclass{ShaderEffect} capability when uploaded
through the \texttt{CNAEXT} \texttt{SetDataRaw} path. Every other renderer still discards that
declaration at the interface default and relies on its known stride/effect combinations;
SDL\_Renderer and the other 2D-only renderers reject vertex-buffer construction outright.
